\documentclass[10pt]{beamer}

% Konfigurasi Tema
\usetheme{Madrid}
\usecolortheme{whale}

% Paket Tambahan
\usepackage[utf8]{inputenc}
\usepackage[T1]{fontenc}
\usepackage{amsmath}
\usepackage{graphicx}
\usepackage{booktabs}
\usepackage{caption}

% Meta Data Presentasi
\title[Optimisasi Portofolio VQE]{Optimisasi Portofolio Berkualitas Tinggi dengan Variational Quantum Eigensolver (VQE)}
\subtitle{Pendekatan Hibrida Menggunakan Weighted CVaR dan CMA-ES}
\author{Nama Presenter} % Ganti dengan nama Anda
\institute[Nama Institusi]{Departemen Fisika / Komputasi \\ Nama Institusi Anda}
\date{\today}

% Penomoran Seksi Otomatis (Slide Pembuka Bagian)
\AtBeginSection[]
{
  \begin{frame}
    \frametitle{Outline Bagian}
    \tableofcontents[currentsection]
  \end{frame}
}

% Memulai Dokumen
\begin{document}

% --- SLIDE 1: JUDUL ---
\begin{frame}
    \titlepage
    \begin{center}
        \textit{\footnotesize [Saran Gambar: Ilustrasi artistik grafik saham digabungkan dengan simbol sirkuit kuantum atau Bola Bloch]}
    \end{center}
\end{frame}

% --- SLIDE 2: DAFTAR ISI ---
\begin{frame}{Daftar Isi Keseluruhan}
    \tableofcontents
\end{frame}

% --- SLIDE 2: LATAR BELAKANG ---
\section{Pendahuluan}
\begin{frame}{Latar Belakang Masalah}
    \begin{itemize}
        \item \textbf{Optimisasi Portofolio:} Masalah fundamental keuangan.
        \begin{itemize}
            \item Tujuan: Alokasi modal untuk \textbf{maksimalkan untung} dan \textbf{minimalkan risiko}.
        \end{itemize}
        \item \textbf{Formulasi QUBO:} Masalah diterjemahkan ke model \textit{Quadratic Unconstrained Binary Optimization}.
        \item \textbf{Tantangan Komputasi:}
        \begin{itemize}
            \item Termasuk kategori \textbf{NP-hard}.
            \item Komputer klasik kesulitan seiring bertambahnya jumlah aset.
        \end{itemize}
        \item \textbf{Solusi Kuantum:} Potensi efisiensi komputasi untuk ruang pencarian yang besar.
    \end{itemize}
    
    \vspace{0.5cm}
    % Ruang untuk Gambar
    \begin{alertblock}{Saran Visualisasi}
        \centering
        Diagram Alur: Data Saham $\rightarrow$ Rumus Matematika (QUBO) $\rightarrow$ Tembok Penghalang (NP-hard) $\rightarrow$ Solusi (Komputer Kuantum)
    \end{alertblock}
\end{frame}

% --- SLIDE 3: PERANGKAT NISQ ---
\begin{frame}{Perangkat Komputasi Era NISQ}
    \begin{columns}
        \column{0.6\textwidth}
        \begin{itemize}
            \item \textbf{Algoritma Kandidat:}
            \begin{itemize}
                \item VQE (\textit{Variational Quantum Eigensolver})
                \item QAOA (\textit{Quantum Approximate Optimization Algorithm})
            \end{itemize}
            \item \textbf{Sifat Hibrida:} Kerjasama CPU dan QPU.
            \begin{itemize}
                \item \textbf{QPU:} Menghitung energi sistem (rumit).
                \item \textbf{CPU:} Optimasi parameter (klasik).
            \end{itemize}
            \item \textbf{Era NISQ:} Perangkat saat ini masih "bising" (\textit{noisy}), namun VQE dinilai cukup tangguh (\textit{robust}).
        \end{itemize}
        
        \column{0.4\textwidth}
        % Ruang untuk Gambar
        \begin{exampleblock}{Ilustrasi}
            \centering
            Gambar skema loop tertutup antara CPU (Klasik) dan QPU (Quantum).
        \end{exampleblock}
    \end{columns}
\end{frame}

% --- SLIDE 4: TANTANGAN TEKNIS ---
\begin{frame}{Tantangan Teknis VQE pada Keuangan}
    \begin{itemize}
        \item \textbf{Perbedaan Domain:}
        \begin{itemize}
            \item \textit{Kimia:} Molekul kompleks, superposisi rumit.
            \item \textit{Keuangan:} Hamiltonian komutatif, lebih sederhana.
        \end{itemize}
        \item \textbf{Ground State:} Solusi berupa \textit{Computational Basis State} (Deretan biner 0 dan 1).
        \item \textbf{Masalah Nilai Ekspektasi:}
        \begin{itemize}
            \item Menggunakan rata-rata energi biasa kurang efektif.
            \item Dapat mengaburkan solusi optimal yang spesifik (jarum dalam jerami).
        \end{itemize}
    \end{itemize}
\end{frame}

% --- SLIDE 5: SOLUSI DIUSULKAN ---
\begin{frame}{Solusi yang Diusulkan}
    \begin{enumerate}
        \item \textbf{Fungsi Biaya Baru (CVaR):}
        \begin{itemize}
            \item Mengganti rata-rata standar dengan \textit{Conditional Value-at-Risk}.
            \item Fokus pada ekor distribusi energi terendah.
        \end{itemize}
        \vspace{0.3cm}
        \item \textbf{Inovasi: Weighted CVaR (WCVaR):}
        \begin{itemize}
            \item Memberi bobot prioritas pada solusi terbaik (ranking atas).
        \end{itemize}
        \vspace{0.3cm}
        \item \textbf{Pengoptimasi CMA-ES:}
        \begin{itemize}
            \item \textit{Covariance Matrix Adaptation Evolution Strategy}.
            \item Bebas turunan (\textit{derivative-free}) untuk menangani lanskap energi kasar.
        \end{itemize}
    \end{enumerate}
\end{frame}

% --- SLIDE 6: PRAPEMROSESAN DATA ---
\section{Landasan Teori}
\begin{frame}{Prapemrosesan Data Saham}
    \textbf{Input:} Matriks harga aset ($P$) sepanjang waktu ($M$).
    
    \vspace{0.5cm}
    \textbf{Tahapan Transformasi:}
    \begin{enumerate}
        \item Hitung \textbf{Return Harian} ($r_{ki}$): Standarisasi perubahan harga.
        \item Hitung \textbf{Ekspektasi Return} ($\mu_i$):
        \begin{itemize}
            \item Merepresentasikan potensi \textcolor{blue}{Keuntungan}.
        \end{itemize}
        \item Hitung \textbf{Matriks Kovariansi} ($\sigma_{ij}$):
        \begin{itemize}
            \item Merepresentasikan tingkat \textcolor{red}{Risiko} dan korelasi antar aset.
        \end{itemize}
    \end{enumerate}

    \vspace{0.5cm}
    \begin{block}{Saran Gambar}
        \centering
        Tabel ilustrasi: Harga Rupiah $\rightarrow$ Persentase Return $\rightarrow$ Vektor $\mu$ \& Matriks $\sigma$.
    \end{block}
\end{frame}

% --- SLIDE 7: MODEL MATEMATIKA ---
\begin{frame}{Model Matematika Portofolio}
    \begin{itemize}
        \item \textbf{Tujuan 1 (Max):} Keuntungan Portofolio ($C_1$).
        \item \textbf{Tujuan 2 (Min):} Risiko Portofolio ($C_2$).
        \item \textbf{Kendala:} Total investasi $\le$ Modal ($B$).
    \end{itemize}
    
    \vspace{0.5cm}
    \textbf{Fungsi Objektif Gabungan:}
    \begin{equation*}
        C' = -\lambda C_1 + (1-\lambda)C_2
    \end{equation*}
    \begin{small}
    Dimana $\lambda$ (0 s.d 1) adalah faktor pengatur selera risiko investor.
    \end{small}

    \vspace{0.5cm}
    \begin{exampleblock}{Saran Gambar}
        \centering
        Grafik \textit{Efficient Frontier} (Kurva Pareto): Sumbu X (Risiko) vs Sumbu Y (Return).
    \end{exampleblock}
\end{frame}

% --- SLIDE 8: QUBO \& PENALTI ---
\begin{frame}{Formulasi QUBO \& Metode Penalti}
    \begin{itemize}
        \item \textbf{Variabel Biner ($x_i$):}
        \begin{itemize}
            \item $1$: Aset dipilih (Beli).
            \item $0$: Aset tidak dipilih.
        \end{itemize}
        \item \textbf{Constraint:} Harus memilih tepat $N/2$ aset.
        \item \textbf{Metode Penalti:}
        \begin{itemize}
            \item Menambahkan suku kuadratik ke fungsi tujuan:
            \[ P = p \left( \sum x_i - \frac{N}{2} \right)^2 \]
            \item Menciptakan "tembok energi" yang tinggi jika aturan dilanggar.
        \end{itemize}
    \end{itemize}

    \vspace{0.3cm}
    \begin{alertblock}{Analogi Fisika}
        \centering
        Mirip energi potensial pegas ($U \sim kx^2$). Memaksa sistem kembali ke posisi setimbang ($N/2$).
        \\ \footnotesize{[Saran Gambar: Grafik parabola "sumur energi"]}
    \end{alertblock}
\end{frame}

% --- SLIDE 9: IMPLEMENTASI VQE ---
\begin{frame}{Implementasi VQE: Dari Keuangan ke Fisika}
    \textbf{1. Mapping Variabel:}
    \[ x_i \in \{0,1\} \xrightarrow{\text{Transformasi}} s_i \in \{+1,-1\} \xrightarrow{\text{Operator}} \hat{Z}_i \]
    
    \vspace{0.3cm}
    \textbf{2. Hamiltonian Sistem:}
    \[ \hat{H} = -\sum h_i \hat{Z}_i - \sum J_{ij} \hat{Z}_i \hat{Z}_j \]
    \begin{itemize}
        \item $h_i$: Representasi Return saham (bias linear).
        \item $J_{ij}$: Representasi Risiko/Kovariansi (interaksi kuadratik).
    \end{itemize}
    
    \vspace{0.3cm}
    \textbf{3. Tujuan:} 
    Mencari energi terendah ($\min \langle H \rangle$) $\equiv$ Profit Maksimal.
\end{frame}

% --- SLIDE 10: KONSEP CVAR ---
\begin{frame}{Conditional Value-at-Risk (CVaR)}
    \begin{columns}
        \column{0.6\textwidth}
        \begin{itemize}
            \item \textbf{Masalah Rata-rata (Expectation Value):}
            \begin{itemize}
                \item Sering menyesatkan jika solusi optimal jarang muncul (probabilitas rendah).
            \end{itemize}
            \item \textbf{Konsep CVaR:}
            \begin{itemize}
                \item Hanya menghitung rata-rata dari $\alpha\%$ sampel terbaik (energi terendah).
                \item Mengabaikan sampel energi tinggi (solusi buruk).
            \end{itemize}
            \item \textbf{Parameter $\alpha$:} Mengatur keseimbangan antara eksplorasi dan stabilitas.
        \end{itemize}
        
        \column{0.4\textwidth}
        \begin{block}{Visualisasi}
            \centering
            [Saran Gambar: Histogram distribusi energi, dengan area $\alpha\%$ terendah diarsir]
        \end{block}
    \end{columns}
\end{frame}

% --- SLIDE 11: INOVASI WCVAR ---
\section{Desain Algoritma}
\begin{frame}{Inovasi: Weighted CVaR (WCVaR)}
    \textbf{Kritik terhadap CVaR Standar:}
    \begin{itemize}
        \item Memberi bobot sama rata pada semua sampel di area "Top $\alpha$".
        \item Juara 1 dan Juara 10 dianggap sama pentingnya.
    \end{itemize}
    
    \vspace{0.3cm}
    \textbf{Solusi WCVaR:}
    \begin{itemize}
        \item Memberi bobot berbeda ($w_k$) berdasarkan peringkat.
        \item Solusi ranking 1 mendapat bobot \textbf{jauh lebih besar}.
        \item Memberi sinyal kuat ke optimizer untuk mengejar solusi elit.
    \end{itemize}

    \vspace{0.3cm}
    \begin{exampleblock}{Perbandingan Bobot}
        \centering
        [Saran Gambar: Grafik garis datar (CVaR) vs Kurva meluruh eksponensial (WCVaR)]
    \end{exampleblock}
\end{frame}

% --- SLIDE 12: OPTIMIZER CMA-ES ---
\begin{frame}{Optimizer: CMA-ES}
    \textbf{Tantangan WCVaR:}
    \begin{itemize}
        \item Menghasilkan lanskap fungsi biaya yang kasar (\textit{non-smooth}) dan berisik.
        \item Optimizer berbasis gradien sering macet.
    \end{itemize}
    
    \vspace{0.5cm}
    \textbf{Mengapa CMA-ES?}
    \begin{itemize}
        \item \textit{Covariance Matrix Adaptation Evolution Strategy}.
        \item \textbf{Evolutionary:} Berbasis populasi, bukan titik tunggal.
        \item \textbf{Derivative-free:} Tidak butuh turunan matematika.
        \item \textbf{Robust:} Tangguh terhadap jebakan minima lokal.
    \end{itemize}
\end{frame}

% --- SLIDE 13: ANSATZ ---
\begin{frame}{Desain Sirkuit Kuantum (Ansatz)}
    \begin{columns}
        \column{0.5\textwidth}
        \textbf{Desain 1: Two-Local Ansatz}
        \begin{itemize}
            \item Struktur standar.
            \item Lapisan Rotasi ($R_y, R_z$) + Entanglement ($CNOT$) berantai.
        \end{itemize}
        \vspace{1cm}
        \begin{block}{Gambar 1a}
            \centering
            [Diagram Sirkuit Two-Local]
        \end{block}
        
        \column{0.5\textwidth}
        \textbf{Desain 2: Block Ansatz}
        \begin{itemize}
            \item Menggunakan modul blok fungsional.
            \item Menciptakan korelasi yang lebih dalam dan kompleks.
        \end{itemize}
        \vspace{1cm}
        \begin{block}{Gambar 1b}
            \centering
            [Diagram Sirkuit Block Ansatz]
        \end{block}
    \end{columns}
\end{frame}

% --- SLIDE 14: SETUP EKSPERIMEN ---
\section{Hasil Numerik}
\begin{frame}{Setup Eksperimen}
    \begin{itemize}
        \item \textbf{Data Pasar Riil (2024):}
        \begin{itemize}
            \item 6 Saham China (A-shares).
            \item 6 Saham Amerika Serikat (US Equities).
        \end{itemize}
        \item \textbf{Ukuran Sistem:} 12 Aset (12 Qubit).
        \begin{itemize}
            \item Cukup kompleks untuk bukti konsep.
            \item Dapat diverifikasi kebenarannya dengan \textit{brute force} klasik.
        \end{itemize}
        \item \textbf{Metrik Penilaian: Top-10 Hit Rate}
        \begin{itemize}
            \item Iterasi dianggap SUKSES jika solusi benar (\textit{Ground State}) masuk dalam \textbf{10 tebakan probabilitas tertinggi}.
        \end{itemize}
    \end{itemize}
\end{frame}

% --- SLIDE 15: HASIL PERFORMA ---
\begin{frame}{Hasil Performa: WCVaR vs CVaR}
    \begin{columns}
        \column{0.6\textwidth}
        \begin{itemize}
            \item \textbf{Stabilitas WCVaR:}
            \begin{itemize}
                \item Unggul dan stabil di berbagai nilai $\alpha$.
                \item Tidak perlu \textit{tuning} parameter yang rumit.
            \end{itemize}
            \item \textbf{Sensitivitas CVaR:}
            \begin{itemize}
                \item Hanya bekerja baik pada $\alpha$ yang sangat kecil.
                \item Performa turun drastis pada $\alpha$ besar.
            \end{itemize}
            \item \textbf{Optimizer:} CMA-ES konvergen jauh lebih cepat dibanding COBYLA.
        \end{itemize}
        
        \column{0.4\textwidth}
        \begin{alertblock}{Grafik Hasil}
            \centering
            [Saran Gambar: Grafik Probability Success vs Iterasi. Garis CMA-ES+WCVaR paling atas]
        \end{alertblock}
    \end{columns}
\end{frame}

% --- SLIDE 16: EVALUASI PEMBOBOTAN ---
\begin{frame}{Evaluasi Resep Pembobotan}
    \textbf{Kami menguji 4 varian pembobotan:}
    \begin{enumerate}
        \item \textbf{CVaR Standar:} Bobot rata (kurang optimal).
        \item \textbf{Energy-based:} Berdasarkan nilai energi (Paling buruk, tidak stabil).
        \item \textbf{Rank-based:} Berdasarkan urutan ranking (Sangat efektif).
        \item \textbf{Piecewise Exponential:} Kustom 3 zona (Terbaik, tapi rumit).
    \end{enumerate}
    
    \vspace{0.5cm}
    \begin{block}{Rekomendasi}
        Metode \textbf{Rank-based (W3)} direkomendasikan karena memberikan keseimbangan terbaik antara kesederhanaan dan performa tinggi.
    \end{block}
\end{frame}

% --- SLIDE 17: KESIMPULAN ---
\section{Kesimpulan}
\begin{frame}{Kesimpulan}
    \begin{itemize}
        \item \textbf{Sinergi adalah Kunci:}
        \begin{itemize}
            \item Penggabungan fungsi biaya \textbf{WCVaR} dan optimizer \textbf{CMA-ES} memberikan hasil paling konsisten.
        \end{itemize}
        \item \textbf{Mengatasi Kendala NISQ:}
        \begin{itemize}
            \item Pendekatan ini efektif menangani lanskap energi yang kasar dan gangguan (\textit{noise}).
        \end{itemize}
        \item \textbf{Fleksibilitas:}
        \begin{itemize}
            \item Jenis Ansatz (Two-local vs Block) tidak memberikan perbedaan performa yang signifikan, memberikan kebebasan desain.
        \end{itemize}
    \end{itemize}
    
    \vspace{0.5cm}
    \centering
    \textbf{Metode ini menjanjikan untuk penerapan optimisasi portofolio praktis di era komputasi kuantum saat ini.}
\end{frame}

% --- SLIDE 18: PENUTUP ---
\begin{frame}
    \centering
    \Huge \textbf{Terima Kasih}
    
    \vspace{1cm}
    \Large Sesi Tanya Jawab (Q\&A)
\end{frame}

\end{document}